\documentclass[11pt,a4paper]{article}

% --- 核心宏包与中文支持 ---
\usepackage[utf8]{inputenc}
\usepackage{xeCJK}
\usepackage{amsmath, amssymb, amsfonts}
\usepackage{listings}
\usepackage{xcolor}
\usepackage{geometry}
\usepackage{tcolorbox}
\usepackage{booktabs}
\usepackage{tabularx}
\usepackage{setspace}
\usepackage{enumitem}
\usepackage[hidelinks]{hyperref}
\usepackage{bookmark}
\usepackage{fancyhdr}

% --- PDF 书签与元数据 ---
\hypersetup{
    colorlinks=false,
    pdfborder={0 0 0},
    bookmarks=true,
    bookmarksnumbered=true,
    bookmarksopen=true,
    bookmarksopenlevel=2,
    pdfcreator={XeLaTeX},
    pdftitle={06 - File I/O & Error Handling},
    pdfauthor={黄子扬},
}

% --- 页面美化设置 ---
\geometry{top=0.7in,bottom=0.8in,left=0.6in,right=0.6in}
\onehalfspacing
\tcbuselibrary{skins,breakable}
\setlength{\parindent}{0pt}
\setlength{\parskip}{0.6em}

% --- 代码样式配置 ---
\definecolor{mainblue}{RGB}{0,51,102}
\definecolor{examred}{RGB}{180,0,0}
\definecolor{codebg}{RGB}{248,248,248}
\definecolor{commentgreen}{RGB}{0,128,0}

\lstset{
    backgroundcolor=\color{codebg},
    basicstyle=\ttfamily\footnotesize,
    keywordstyle=\color{blue}\bfseries,
    commentstyle=\color{commentgreen},
    stringstyle=\color{orange},
    breaklines=true,
    showstringspaces=false,
    frame=leftline,
    xleftmargin=2em,
    numbers=left,
    numbersep=5pt,
    numberstyle=\tiny\color{gray},
    language=Python,
    morekeywords={None, True, False, with, as, try, except, raise, import, def, for, in, if, else}
}

% --- 自定义教学框 ---
\newtcolorbox{concept}[1]{colback=blue!5, colframe=mainblue, fonttitle=\bfseries, title={{【Concept / 核心概念】#1}}, arc=3pt, breakable}
\newtcolorbox{warning}[1]{colback=red!5, colframe=examred, fonttitle=\bfseries, title={{【Warning / 避坑指南】#1}}, arc=3pt, breakable}
\newtcolorbox{examplebox}[1]{colback=green!2, colframe=green!40!black, fonttitle=\bfseries, title={{【Example / 实战案例】#1}}, arc=2pt, breakable}

% --- 页眉页脚 ---
\pagestyle{fancy}
\fancyhf{}
\fancyhead[L]{\small\bfseries 06 - File I/O \& Error Handling / 文件IO与异常处理}
\fancyhead[R]{\small\thepage}
\renewcommand{\headrulewidth}{0.5pt}

% --- 文档信息 ---
\title{\Huge \bfseries 06 - File I/O \& Error Handling}
\author{\Large \textbf{哈尔滨工业大学 \quad 黄子扬}}
\date{2026年6月8日}

\begin{document}

\maketitle
\tableofcontents
\newpage

% ======================================================================
\section{File Systems \& Paths / 文件系统与路径}
% ======================================================================

\begin{concept}{Hierarchy and Paths / 层级与路径}
    File systems organize files in a hierarchy of \textbf{directories} (folders). At the top is the \textbf{root directory}. \\
    文件系统将文件组织在目录（文件夹）的层级结构中。最顶层是根目录。

    Each file in a directory has a file name that must be unique from any other file in the same directory. Directories can contain \textbf{subdirectories} (subfolders).

    A \textbf{file path} is a string representing the location of a file, starting from the root directory. \\
    \textbf{文件路径}是从根目录开始定位文件的字符串。
\end{concept}

\begin{warning}{OS Differences / 操作系统路径差异}
    How you write file paths depends on your operating system:
    \begin{itemize}
        \item \textbf{Windows}: Uses backslashes (\texttt{\textbackslash}), e.g., \texttt{C:\textbackslash Users\textbackslash David\textbackslash Documents\textbackslash hello.txt}.
        \item \textbf{macOS}: Uses forward slashes (\texttt{/}), e.g., \texttt{/Users/David/Documents/hello.txt}.
        \item \textbf{Ubuntu Linux}: Uses forward slashes (\texttt{/}), e.g., \texttt{/home/David/Documents/hello.txt}.
    \end{itemize}
\end{warning}

% ======================================================================
\section{The open() Function / 打开文件}
% ======================================================================

\subsection{Opening Modes / 打开模式}

\begin{concept}{Creating a File Object / 创建文件对象}
    Before you can read or write a file, you must open it using Python's built-in \texttt{open()} function. This function creates a \textbf{file object}. \\
    读写文件之前必须先用 \texttt{open()} 打开，这会创建一个文件对象。

    \texttt{open(file='path/to/file', mode='...')}
    \begin{itemize}[itemsep=0.5em]
        \item \textbf{\texttt{'r'} (Read)}：Read only (default if mode is omitted). File must exist. / 只读。
        \item \textbf{\texttt{'w'} (Write)}：Overwrites existing file or creates a new one. \textbf{WARNING}: erases existing content! / 覆写（会清空原内容）。
        \item \textbf{\texttt{'a'} (Append)}：Adds content to the end of the file. Creates the file if it doesn't exist. / 追加（在末尾添加内容）。
    \end{itemize}
\end{concept}

\subsection{File Attributes / 文件属性}

\begin{examplebox}{Checking File Attributes / 查看文件属性}
\begin{lstlisting}
my_file = open(file='sample.txt', mode='w')

print("Name of the file: ", my_file.name)     # sample.txt
print("Closed or not: ", my_file.closed)       # False
print("Opening mode: ", my_file.mode)          # w
\end{lstlisting}
\end{examplebox}

\subsection{Writing to Files / 写入文件}

\begin{examplebox}{write() Method / write() 方法}
    The \texttt{write()} method writes any string to an open file. \textbf{IMPORTANT}: It does \textbf{NOT} add a newline character (\texttt{\textbackslash n}) automatically — you must add it manually! \\
    \texttt{write()} 不会自动添加换行符，必须手动加 \texttt{\textbackslash n}！
\begin{lstlisting}
my_file = open(file='sample.txt', mode='w')
my_file.write('Hello file!')
my_file.write('\n')                 # Manually add newline
my_file.write('This is me!')
my_file.close()                     # Always close to flush data!
\end{lstlisting}
\end{examplebox}

\subsection{The close() Method / 关闭文件}

\begin{concept}{Why close()? / 为什么必须 close()？}
    The \texttt{close()} method flushes any unwritten information and closes the file object, after which no more writing can be done. \\
    \texttt{close()} 刷写缓冲区中的未保存数据并关闭文件。不关闭可能导致数据丢失！
\end{concept}

% ======================================================================
\section{Reading Files / 读取文件}
% ======================================================================

\begin{concept}{Four Ways to Read / 四种读取方式}
    \begin{itemize}[itemsep=0.8em]
        \item \textbf{\texttt{.read()}}：Reads the \textbf{entire} file as one single string. / 整个文件读入一个字符串。
        \item \textbf{\texttt{.readline()}}：Reads just \textbf{one line} (up to and including the newline character). / 仅读取一行。
        \item \textbf{\texttt{.readlines()}}：Reads the whole file and returns a \textbf{list of strings} — each line is one element. / 返回字符串列表，每个元素是一行。
        \item \textbf{\texttt{for line in file}}：Iterates line by line — the most \textbf{memory-efficient} way! / 逐行遍历，最省内存。
    \end{itemize}
\end{concept}

\begin{examplebox}{Reading Methods Comparison / 四种读取方式对比}
\begin{lstlisting}
# Method 1: .read() — entire file as one string
my_file = open(file='sample.txt', mode='r')
file_output = my_file.read()
my_file.close()
print(file_output)

# Method 2: .readline() — just one line
my_file = open(file='sample.txt', mode='r')
file_output = my_file.readline()    # Reads first line
my_file.close()

# Method 3: .readlines() — list of lines
my_file = open(file='sample.txt', mode='r')
file_output = my_file.readlines()
my_file.close()
print(file_output)                  # ['Hello file!\n', 'This is me!\n']

# Method 4: for line in file — BEST for large files!
my_file = open(file='sample.txt', mode='r')
for line in my_file:                # Most memory-efficient
    print(line)
my_file.close()
\end{lstlisting}
\end{examplebox}

% ======================================================================
\section{Working with CSV Files / CSV 文件处理}
% ======================================================================

\begin{concept}{What is CSV? / 什么是 CSV？}
    CSV stands for \textbf{Comma Separated Values}. It is a simple but powerful way to store data in a table format. \\
    CSV 是一种用逗号分隔值的简单表格数据格式。
\end{concept}

\subsection{Reading CSV / 读取 CSV}

\begin{examplebox}{Using csv.reader() and DictReader / CSV 读取器}
\begin{lstlisting}
import csv

# Method 1: csv.reader — returns lists
csv_file = open("sample.csv", "r")
reader = csv.reader(csv_file, skipinitialspace=True)
for row in reader:
    print(row)                      # ['1', 'Rose', '30', 'Grenoble']
csv_file.close()

# Method 2: csv.DictReader — returns dicts (RECOMMENDED!)
csv_file = open("sample.csv", "r")
reader = csv.DictReader(csv_file)
for row in reader:
    print(row)
    # {'id': '1', 'name': 'Rose', 'age': '30', 'city': 'Grenoble'}
csv_file.close()
\end{lstlisting}
\end{examplebox}

\subsection{Writing CSV / 写入 CSV}

\begin{examplebox}{Using csv.writer() / CSV 写入器}
\begin{lstlisting}
# Write CSV — use "w" to overwrite, "a" to append
csv_file = open("sample.csv", "w")
writer = csv.writer(csv_file)
writer.writerow(['5', 'Rose', '30', 'Grenoble'])
writer.writerow(['6', 'Nona', '23', 'Toulouse'])
csv_file.close()
\end{lstlisting}
\end{examplebox}

\begin{warning}{'w' vs 'a' for CSV / CSV 写入模式选择}
    If you do NOT want to overwrite your existing file, use \textbf{\texttt{'a'}} (append) instead of \textbf{\texttt{'w'}} (write). \\
    用 \texttt{'w'} 会清空原文件内容，用 \texttt{'a'} 会在末尾追加！
\end{warning}

% ======================================================================
\section{Error Handling \& Debugging / 异常处理与调试}
% ======================================================================

\subsection{Error Types / 错误类型}

\begin{concept}{Errors vs Exceptions / 语法错误 vs 运行时异常}
    \begin{itemize}[itemsep=0.8em]
        \item \textbf{SyntaxError / 语法错误}：Caught by the IDE \textbf{before} execution. The script cannot run at all. Examples: forgetting a colon \texttt{:}, forgetting to close a bracket, invalid variable names.
        \item \textbf{Exceptions / 运行时异常}：Occur \textbf{during} execution. Examples: \texttt{IndexError} (list index out of range), \texttt{NameError} (variable not defined), \texttt{TypeError} (wrong type operation).
    \end{itemize}
\end{concept}

\begin{examplebox}{Common Error Examples / 常见错误类型}
\begin{lstlisting}
# SyntaxError: invalid name (starts with digit)
# 3_apples = ["apple", "apple", "apple"]

# SyntaxError: forgot to close brackets
# three_apples = ["apple", "apple", "apple"

# SyntaxError: forgot colon and wrong indentation
# for i in range(3)
# print(three_apples[i])

# IndexError: index out of range
three_apples = ['A', 'B', 'C']
# print(three_apples[4])         # IndexError!

# TypeError: cannot concatenate str + int
# print("hello " + 6)            # TypeError!

# ValueError: invalid literal for float()
converting_data_types = "I am a string"
# print(float(converting_data_types))  # ValueError!
\end{lstlisting}
\end{examplebox}

\subsection{The Try-Except Block / 捕获异常}

\begin{concept}{try-except Logic / try-except 逻辑}
    \texttt{try-except} acts like \texttt{if-else} but for exceptions. When the code in the \texttt{try} block fails (raises an exception), the \texttt{except} block runs instead. \\
    类似于 \texttt{if-else}，但专门处理异常。\texttt{try} 块出错时自动执行 \texttt{except} 块。
\end{concept}

\begin{examplebox}{Handling User Input / 处理用户输入案例}
\begin{lstlisting}
user_input = input("Tell me your name please: ")

try:
    user_input = float(user_input)
    print("That is a number, not your name!")
except:
    print("Thank you,", user_input)
\end{lstlisting}
    If the user enters a string that can't be converted to float, the \texttt{except} block runs. If the conversion succeeds, the \texttt{try} block completes.
\end{examplebox}

\subsection{Debugging Tips / 调试技巧}

\begin{concept}{Using print() for Debugging / 用 print() 调试}
    Use the \texttt{print()} function to keep track of your variables and outputs. This helps you catch the point where your code misbehaves. \\
    在关键位置插入 \texttt{print()} 来追踪变量值，定位代码出错的位置。
\end{concept}

% ======================================================================
\section{PEP-8 Style Guide / PEP-8 编程规范}
% ======================================================================

\begin{concept}{Professional Coding / 像专家一样写代码}
    Writing nice, easy-to-understand code is crucial. For Python, the official style guide is \textbf{PEP-8}. \\
    写出易读、易懂的代码至关重要。Python 的官方代码风格指南是 PEP-8。

    Conventions are specified in a strict manner, but remember: don't spend more time with the style guide than with your actual code.
\end{concept}

\subsection{Comments and Docstrings / 注释与文档字符串}

\begin{warning}{DOs and DO NOTs for Comments / 注释的规范}
    \textbf{DOs / 应该做的}：
    \begin{itemize}
        \item Add explanatory comments before lines that are not entirely obvious.
        \item Use the same indentation as the code they describe.
        \item Add docstrings as the first line of every function and module.
    \end{itemize}

    \textbf{DO NOTs / 不应该做的}：
    \begin{itemize}
        \item Do NOT add self-explanatory comments: \texttt{x = x + 1  \# increment x} is BAD.
        \item Do NOT add comments that are wrong — worse than no comment at all!
    \end{itemize}
\end{warning}

\subsection{Whitespaces / 空格规范}

\begin{concept}{Where to Put Spaces / 空格的位置}
    \textbf{DO put spaces / 应该加空格的}：
    \begin{itemize}
        \item After commas, colons, and semicolons.
        \item Around \texttt{=}, \texttt{+=} in assignments.
        \item Around most binary operators, e.g., \texttt{2 + 2}.
    \end{itemize}

    \textbf{DO NOT put spaces / 不应该加空格的}：
    \begin{itemize}
        \item Between a function name and \texttt{()}: \texttt{func(x)} NOT \texttt{func (x)}.
        \item Around \texttt{=} in keyword arguments: \texttt{n=42} NOT \texttt{n = 42}.
        \item After \texttt{(} and before \texttt{)}.
    \end{itemize}
\end{concept}

\begin{examplebox}{Good vs Bad Formatting / 好代码 vs 坏代码}
\begin{lstlisting}
# Bad / 坏代码
def my_function ( a,b,c ) :
    some_dict={"a" : "b", 4:5}
    another_function(a+b+c, n = 42)

# Good / 好代码
def my_function(a, b, c):
    some_dict = {"a": "b", 4: 5}
    another_function(a + b + c, n=42)
\end{lstlisting}
\end{examplebox}

\subsection{Blank Lines / 空行规范}

\begin{concept}{How to Use Blank Lines / 如何使用空行}
    \begin{itemize}[itemsep=0.5em]
        \item Separate function definitions by \textbf{two} blank lines.
        \item Separate code blocks (loops, if-statements) by a \textbf{one} blank line.
        \item Separate logical segments in your code by line breaks.
        \item \textbf{Rule of thumb}：No more than 4 lines without a blank line.
        \item Break up long lines of code into multiple ones — either logically (intermediate variables) or syntactically (line continuation with \texttt{\textbackslash}).
    \end{itemize}
\end{concept}

\begin{examplebox}{Line Continuation / 行续接}
\begin{lstlisting}
# Explicit continuation with backslash
total_income = gross_income \
               - taxes \
               - student_loan_interest
\end{lstlisting}
\end{examplebox}

\subsection{Naming Conventions / 命名约定}

\begin{concept}{Standard Naming Patterns / 标准命名模式}
    \begin{itemize}[itemsep=0.5em]
        \item \textbf{Variables and functions}：\texttt{lower\_case\_with\_underscores} (snake\_case) — e.g., \texttt{translation\_dict}, \texttt{word\_counter}.
        \item \textbf{Constants}：\texttt{UPPER\_CASE\_WITH\_UNDERSCORES} — e.g., \texttt{MAX\_SIZE}.
        \item \textbf{Classes}：\texttt{UpperCamelCase} — e.g., \texttt{HREmployee}.
        \item \textbf{Avoid built-in names}：Do NOT use \texttt{list}, \texttt{str}, \texttt{int} as variable names.
    \end{itemize}
\end{concept}

% ======================================================================
\section{Quick Reference / 速查表}
% ======================================================================

\begin{table}[h]
\centering
\small
\renewcommand{\arraystretch}{1.3}
\begin{tabular}{@{}ll@{}}
\toprule
\textbf{Operation / 操作} & \textbf{Syntax / 语法} \\ \midrule
打开文件(读) & \texttt{open('f.txt', 'r')} \\
打开文件(写) & \texttt{open('f.txt', 'w')} \\
打开文件(追加) & \texttt{open('f.txt', 'a')} \\
写入内容 & \texttt{f.write(str)} \\
读取全部 & \texttt{f.read()} \\
读取一行 & \texttt{f.readline()} \\
读取所有行 & \texttt{f.readlines()} \\
逐行遍历 & \texttt{for line in f:} \\
关闭文件 & \texttt{f.close()} \\
CSV 读(Dict) & \texttt{csv.DictReader(f)} \\
CSV 写 & \texttt{csv.writer(f).writerow(row)} \\
捕获异常 & \texttt{try: ... except: ...} \\
命名规范(变量/函数) & \texttt{snake\_case} \\
命名规范(常量) & \texttt{UPPER\_CASE} \\
命名规范(类) & \texttt{CamelCase} \\ \bottomrule
\end{tabular}
\end{table}

\end{document}
